\begin{frame}[fragile]{Projektstruktur}
\small
\begin{tabularx}{\textwidth}{>{\ttfamily}l X}
\toprule
datasets/wider\_face/ & Rohdaten und Splits des WIDER-FACE-Datensatzes. \\
annotations/ & COCO-Annotationen fuer Torchvision-Training und Evaluation. \\
trained\_models/ & Trainierte Gewichte und YOLO-Run-Artefakte. \\
model\_results/ & CSV/JSON-Ergebnisse, Plots und Trainingshistorien. \\
Videos/lab\_outputs/ & Geblurrte Videos, Previews und selektive Testlaeufe. \\
\bottomrule
\end{tabularx}
\vspace{0.8em}
\begin{itemize}
    \item Die Trennung macht Training, Evaluation und Videoausgabe reproduzierbar.
    \item Notebooks dokumentieren die Schritte; Skripte kapseln die ausfuehrbaren Laeufe.
\end{itemize}
\end{frame}
